\section{A sharper quantitative certificate}\label{sec:refinement}

The inherited weighted construction used a Markov bound for a smooth
polynomial's degree. Its explicit search cutoff was
\[
 R_{\mathrm{old}}=480+2{,}000{,}000\,s_q,\qquad
 s_q=|\mathcal P|,\quad \mathcal P=\{P:\deg P\le480\}.
\]
The new refinement uses the variance of that degree. It keeps the
74-point witness and the same proof of progression avoidance.

\subsection{A general concentration lemma}

\begin{lemma}[Positive coefficient from a concentrated generating function]
\label{lem:concentration}
Let \(H(u)=\sum_{n\ge0}H_nu^n\) have nonnegative coefficients and
\(0<H(z)<\infty\) for \(0<z<1\). Define a random variable \(X\) by
\[
 \Prob(X=n)=\frac{H_nz^n}{H(z)}.
\]
Suppose \(J(u)=P(u)-Q(u)\), where \(P,Q\) have nonnegative coefficients,
all exponents in \(P\) are at most \(L\), and
\[
 J(z)>\gamma>0,\qquad P(z)\le C.
\]
If \(\E X\le U\), \(\Var X\le V\), and \(t>0\) satisfies
\(CV/t^2<\gamma\), then for every integer \(R\ge L+U+t\),
\[
 \sum_{n=0}^R[u^n](H(u)J(u))\frac{z^n}{H(z)}>0.
\]
In particular, some coefficient through \(R\) is positive. If the
coefficients are integers, some coefficient is at least one.
\end{lemma}

\begin{proof}
The positive contribution omitted beyond \(R\) is
\[
 \sum_{d\le L}[u^d]P(u)\,z^d\,\Prob(X>R-d)
 \le P(z)\Prob(X>U+t)
 \le \frac{CV}{t^2}
\]
by Chebyshev's inequality, since \(\E X\le U\).
Omitting the negative tail can only increase the truncated signed sum.
Its normalized value therefore exceeds \(\gamma-CV/t^2>0\).
\end{proof}

This is an elementary concentration argument. The contribution here is
its explicit application and certified numerical effect on the existing
polynomial-ring counterexample, not a claim that Chebyshev's inequality
or coefficient extraction is new.

\subsection{Exact integer formulas for all prime powers}

Take \(z=127/128\) and \(L=480\). Write \(I_q(d)\) for the number of monic
irreducibles of degree \(d\) over \(\F_q\), computed recursively by
\begin{equation}\label{eq:irreducibles}
 dI_q(d)=q^d-\sum_{\substack{e\mid d\\e<d}}eI_q(e).
\end{equation}
For the smooth generating function in
Lemma~\ref{lem:replacement}, the exponents of the different irreducibles
are independent geometric random variables with ratios \(z^d\).
Consequently
\begin{align}
 \mu_q:=\E X
   &=\sum_{d=1}^{480}I_q(d)\frac{d z^d}{1-z^d},\\
 v_q:=\Var X
   &=\sum_{d=1}^{480}I_q(d)\frac{d^2 z^d}{(1-z^d)^2}.
\end{align}
No numerical probability calculation is necessary. Define the integers
\begin{align}
 U_q&=\sum_{d=1}^{480}
 \left\lceil\frac{I_q(d)d\,127^d}{128^d-127^d}\right\rceil,\\
 V_q&=\sum_{d=1}^{480}
 \left\lceil\frac{I_q(d)d^2\,127^d128^d}
 {(128^d-127^d)^2}\right\rceil,\\
 R_q^*&=480+U_q+\left\lceil\sqrt{7400V_q}\right\rceil.
 \label{eq:Rstar}
\end{align}
Termwise ceilings give \(\mu_q\le U_q\) and \(v_q\le V_q\).

\begin{theorem}[Improved finite search cutoff]\label{th:cutoff}
For the 74-point witness, an integer \(R\le R_q^*\) satisfies
\([u^R]H(u)J(u)\ge1\). In particular
\begin{equation}\label{eq:explicitgap}
 \beta_q\ge m_q+(q-1)c_q q^{-R_q^*-1}>m_q.
\end{equation}
\end{theorem}

\begin{proof}
Use \(P(u)=\sum_{b\in B}u^b\) and
\(Q(u)=\sum_{a\in A}u^a\).
Equation~\eqref{eq:Jpositive} gives \(\gamma=3/200\), and
\(P(z)\le74\). Put \(t=\lceil\sqrt{7400V_q}\rceil\).
Then
\[
 \frac{74V_q}{t^2}\le\frac{74}{7400}=\frac1{100}.
\]
The normalized truncated sum is greater than
\[
 \frac3{200}-\frac1{100}=\frac1{200}>0.
\]
Lemma~\ref{lem:concentration} provides the positive integer coefficient.
Equation~\eqref{eq:delta} and \(R\le R_q^*\) prove
\eqref{eq:explicitgap}.
\end{proof}

\subsection{The binary case}

Exact integer arithmetic gives the following value. It is shown in four
concatenated blocks to preserve legibility:
\[
\begin{array}{l}
R_2^*={}1493233132307651501440182722716026917359\\[-2pt]
\phantom{R_2^*={}}6595711329751227043997719117349462629005\\[-2pt]
\phantom{R_2^*={}}6106035344506094657698429655284082117452\\[-2pt]
\phantom{R_2^*={}}824668452354378286162187 .
\end{array}
\]
The line breaks are concatenations, not sums. The machine-readable exact
integer is also in \path{data/density_certificate.json}.
Its bit length is 476, and the stronger inequality
\begin{equation}\label{eq:Rbinary}
 R_2^*+13<2^{476}
\end{equation}
was checked exactly. The inherited cutoff is about \(174{,}581.648\)
times larger.

To obtain a simple rigorous lower bound for \(c_2\), use
\(I_2(d)\le2^d/d\) and \(\log(1+t)\le t\):
\[
 c_2=\frac12\prod_{d=1}^{480}(1+2^{-d})^{-I_2(d)}
 \ge\frac12\exp\left(-\sum_{d=1}^{480}\frac1d\right)
 \ge\frac1{960e}>\frac1{2880}>2^{-12}.
\]
Here \(\sum_{d\le480}1/d\le1+\log480\) and \(e<3\).

\begin{corollary}[Explicit strict margin]\label{cor:binary}
\[
 \boxed{\displaystyle \beta_2>m_2+2^{-2^{476}}.}
\]
More sharply, \(\beta_2>m_2+2^{-R_2^*-13}\).
\end{corollary}

\begin{proof}
The additional density in \eqref{eq:explicitgap} is
\(c_2\,2^{-R_2^*-1}>2^{-R_2^*-13}\).
Now use \eqref{eq:Rbinary}.
\end{proof}

\begin{table}[htbp]
\centering
\begin{tabular}{@{}rrrr@{}}
\toprule
\(q\) & \(R_q^*<2^b\) & \(R_{\mathrm{old}}<2^{b_{\mathrm{old}}}\)
 & \(R_{\mathrm{old}}/R_q^*\)\\
\midrule
2 & \(b=476\) & \(b_{\mathrm{old}}=494\) & \(174{,}581.65\)\\
3 & \(756\) & \(774\) & \(175{,}107.59\)\\
4 & \(956\) & \(973\) & \(175{,}282.72\)\\
5 & \(1110\) & \(1127\) & \(175{,}370.25\)\\
7 & \(1343\) & \(1360\) & \(175{,}457.75\)\\
\bottomrule
\end{tabular}
\caption{Certified integer cutoffs. The last column is rounded for display;
all underlying cutoff calculations use integers.}
\end{table}

The new margin remains extraordinarily small. The proof does not exhibit
the first positive coefficient \(R\), because searching up to this cutoff
is infeasible. It proves that such an \(R\) exists within an explicit range.
Thus the improvement is a stronger quantitative theorem, not a practical
enumeration of an extremal polynomial set.
